\documentclass[11pt, oneside]{article}   	
\usepackage{geometry}                		
\geometry{letterpaper}                   		
\usepackage[parfill]{parskip}    		
\usepackage{graphicx}						
\usepackage{amssymb, enumerate}

%begin document
\begin{document}
\begin{center} {\Large{MATH F113X: Homework 3}} \end{center}
\vspace{.3in}

\fbox{\parbox{\textwidth}{
\begin{itemize}
\item Answer the following problems from the Fair Division section (p. 108):
\begin{quote}\# A, 4, 6, B, 8, C \end{quote}
\end{itemize}
}}
\vfill

{\bf Problem A:}
Ahmed, Emily, Tiana, and Levi buy a large pizza for \$24. It is half pepperoni and half veggie and cut into 8 slices.
\begin{enumerate}
	\item[a.] Suppose Ahmed doesn't care whether the slice is pepperoni or veggie. What is the value a slice of veggie for Ahmed?
	\item[b.] Suppose Emily doesn't eat meat of any kind. What is the value a slice of veggie for Emily? What is the value a slice of pepperoni for Emily?
	\item[c.] Suppose Tiana values a slice of pepperoni twice as much as a slice of veggie. What is the value a slice of veggie for Tiana? What is the value a slice of pepperoni for Tiana?
	\item[d.] Are two slices of veggie a fair share for Tiana? Justify your answer.
	\item[e.] Are two slices of veggie a fair share for Emily? Justify your answer.
\end{enumerate}
\vfill

\textbf{Problem B:} Two people, Aeden (A) and Beyonc\'{e} (B), are splitting a box of 6 labubu toys from the Macaron Series worth \$400. The chart below shows how each person values each toy, given in dollars.

\begin{tabular}{| c || c | c | c | c | c | c |}
\hline
& \multicolumn{6}{c}{Labubu Toy}\\
& Soymilk (S)&Lychee (L)&Grape (G) &Coconut (C) &Toffee (T) & Sesame Bean (SB) \\
\hline
Aeden &100&100&50&50&50&50\\
Beyonc\'{e}&100&80&80&80&40&20\\
\hline
\end{tabular}
	\begin{enumerate}
	\item Divide the box of toys using Divider-Chooser with Aeden as the Divider and Beyonc\'{e} as the Chooser. Indicate the final assignment of toys to people and the relative values of each share.
	\item Divide the box of toys using Divider-Chooser with Beyonc\'{e}  as the Divider and Aeden as the Chooser. Indicate the final assignment of toys to people and the relative values of each share.

	\end{enumerate}

        \newpage
	
\textbf{Problem C:} Three people, Aimee (A), Benito (B), and Cicely (C), are dividing a piece of land using the lone-divider method. The values of the three pieces of land in the eyes of each player are shown below.

\begin{tabular}{| c || c | c | c |}
\hline
& piece 1 & piece 2 & piece 3 \\
A&30\%&30\% &40\%\\
B&25\%&25\%&50\%\\
C &33.3\%&33.3\%&33.3\%\\
\hline
\end{tabular}

	\begin{enumerate}
	\item Who was the divider and how can you tell?
	\item Explain what happens in the next step of the lone-divider method and why?
	\item Suppose Cicely takes piece 1. What is the value of remaining land to Aimee and Benito?
	\item Suppose Aimee and Benito split the combined land using the Divider-Chooser method what is the minimum share of the original land each would receive? Would it be fair?
	\end{enumerate}
\vfill
Remember to write up your homework solutions according to the homework writeup guidelines. Confirm your work is correct or fix your errors before you submit it. \\
\vfill
\textbf{Rubric:}

\begin{description}
\item[2 points:] You provided a complete answer, with supporting work, written up clearly
\item[1 point:] Some attempt at a solution, but incomplete writeup / unclear / illegible / no answer
\item[1 point:] Only an answer, with no supporting work 
\item[0 points:] Missing.
\end{description}



\end{document}  
%%% Local Variables:
%%% mode: LaTeX
%%% TeX-master: t
%%% End:
